\documentclass[preview]{standalone}

\usepackage{amsmath}
\usepackage{amssymb}
\usepackage{stellar}
\usepackage{definitions}

\begin{document}

\id{euclidean-domains-exercises}
\genpage

\section{Exercises}

\begin{snippetexercise}{euclidean-quadratic-rings-exercise}{Euclidean quadratic rings}
    Study whether the quadratic rings
    \[
        \integers[\sqrt d]
    \]
    are Euclidean and, when \(d\equiv1\pmod4\), do the same for
    \[
        \integers\left[\frac{1+\sqrt d}{2}\right].
    \]
    In particular, compare the familiar Euclidean examples
    \[
        \integers[i],
        \qquad
        \integers[\sqrt{-2}],
        \qquad
        \integers[\sqrt2],
    \]
    and the ring of integers of \(\rationalnumbers(\sqrt{-3})\),
    \[
        \integers\left[\frac{1+\sqrt{-3}}{2}\right].
    \]
\end{snippetexercise}

\begin{snippetexercise}{gaussian-integers-gcd-bezout-exercise}{GCD and Bézout identity in the Gaussian integers}
    On \(\integers[i]\), consider the Euclidean function
    \[
        S\colon\integers[i]\difference\{0\}\fromto\naturalnumbers,
        \qquad
        S(a+ib)=a^2+b^2.
    \]
    Find a greatest common divisor and a Bézout identity for
    \[
        x=7+17i,
        \qquad
        y=8-14i.
    \]
\end{snippetexercise}

\begin{snippetsolution}{gaussian-integers-gcd-bezout-solution}{GCD and Bézout identity in the Gaussian integers}
    The Euclidean values are
    \[
        S(x)=338,
        \qquad
        S(y)=260.
    \]
    For the first division,
    \[
        \frac{x}{y}
        =
        \frac{(7+17i)(8+14i)}{8^2+14^2}
        =
        \frac{-182+234i}{260}
        =
        -\frac{7}{10}+\frac{117}{130}i.
    \]
    Choose the nearby Gaussian integer \(q=-1+i\). Then
    \begin{align*}
        r
        &=
        x-qy\\
        &=
        7+17i-(-1+i)(8-14i)\\
        &=
        1-5i.
    \end{align*}
    Since \(S(r)=26<260=S(y)\), the first division is
    \[
        x=(-1+i)y+(1-5i).
    \]
    The next division is exact:
    \[
        \frac{y}{1-5i}
        =
        \frac{(8-14i)(1+5i)}{1^2+5^2}
        =
        \frac{78+26i}{26}
        =
        3+i,
    \]
    so
    \[
        y=(3+i)(1-5i).
    \]
    Thus a greatest common divisor, up to a Gaussian unit, is
    \(1-5i\). Reversing the first division gives
    \[
        1-5i=x-(-1+i)y=x+(1-i)y,
    \]
    which is the required Bézout identity.
\end{snippetsolution}

\end{document}
